\begin{kbvrule}
\noindent{\kbvsecfont ಸೂತ್ರ :}\hspace{0.5em}ಸಂಧಿಕಾರ್ಯದಲ್ಲಿ ಪೂರ್ವಶಬ್ದದ ಅಂತ್ಯದಲ್ಲಿರುವ ಕ, ಚ, ಟ, ತ, ಪ ವ್ಯಂಜನಗಳಿಗೆ ಯಾವುದೇ ವರ್ಣ ಪರವಾದರೂ ಕ್ರಮವಾಗಿ ಗ, ಜ, ಡ, ದ, ಬ ವ್ಯಂಜನಗಳು ಆದೇಶಗಳಾಗಿ ಬರುವುದಾದರೆ ಅದು ಜಶ್ತ್ವ ಸಂಧಿ ಯೆನಿಸುವುದು.
\end{kbvrule}

\kbvsechead{(೧)\hspace{0.45em}ಶ್ಚುತ್ವ ಸಂಧಿ :}

\noindent\hspace*{2.2em}ಈ ಉದಾಹರಣೆಗಳನ್ನು ಗಮನಿಸಿ...\par

\indent ({\kbvlatin i}) ಮನಸ್‌ + ಚಂಚಲ = ಮನಶ್‌ + ಚಂಚಲ = ಮನಶ್ಚಂಚಲ (ಸಕಾರಕ್ಕೆ ಚಕಾರ ಪರವಾಗಿದ್ದು ಸಕಾರಕ್ಕೆ ಶಕಾರಾದೇಶವಾಗಿದೆ)\par

\indent ({\kbvlatin ii}) ತಕಾರ + ಚಕಾರ = ತಕಾರ + ಚಕಾರ = ತಕಾರಚಕಾರ (ತಕಾರಕ್ಕೆ ಚಕಾರ ಪರವಾಗಿದ್ದು ತಕಾರಕ್ಕೆ ಚಕಾರಾದೇಶವಾಗಿದೆ)\par

\indent ({\kbvlatin iii}) ಬೃಹತ್‌ + ಛತ್ರ = ಬೃಹಚ್‌ + ಛತ್ರ = ಬೃಹಚ್ಛತ್ರ (ತಕಾರಕ್ಕೆ ಛಕಾರ ಪರವಾಗಿದ್ದು ತಕಾರಕ್ಕೆ ಚಕಾರಾದೇಶವಾಗಿದೆ)\par

\indent ({\kbvlatin iv}) ಮನಸ್‌ + ಶುದ್ಧಿ = ಮನಶ್‌ + ಶುದ್ಧಿ = ಮನಶ್ಶುದ್ಧಿ (ಸಕಾರಕ್ಕೆ ಶಕಾರ ಪರವಾಗಿದ್ದು ಸಕಾರಕ್ಕೆ ಶಕಾರಾದೇಶವಾಗಿದೆ)\par

\indent —ಇವು ಶ್ಚುತ್ವ ಸಂಧಿಗೆ ಉದಾಹರಣೆಗಳು. ಶ್ಚುತ್ವ ಸಂಧಿ ಯಾಗುವ ಸಂದರ್ಭಗಳನ್ನು ಕೆಳಗೆ ಕೊಟ್ಟಿರುವ ಸೂತ್ರದಿಂದ ತಿಳಿಯಬಹುದು.\par

\begin{kbvrule}
\noindent{\kbvsecfont ಸೂತ್ರ :}\hspace{0.5em}ಪೂರ್ವ ಪದದ ಅಂತ್ಯದಲ್ಲಿ ಸಕಾರವಾಗಲೀ ತವರ್ಗದಲ್ಲಿನ ಯಾವುದಾದರೊಂದು ವರ್ಣವಾಗಲೀ ಇದ್ದು, ಆ ವರ್ಣಗಳಿಗೆ ಶಕಾರವಾಗಲೀ ಚವರ್ಗದಲ್ಲಿನ ಯಾವುದಾದರೊಂದು ವರ್ಣವಾಗಲೀ ಪರವಾದಾಗ, ಸಕಾರಕ್ಕೆ ಶಕಾರವೂ ತವರ್ಗಕ್ಕೆ ಚವರ್ಗವೂ ಆದೇಶಗಳಾಗಿ ಬರುವುದಾದರೆ ಅದು ಶ್ಚುತ್ವ ಸಂಧಿಯೆನಿಸುವುದು.
\end{kbvrule}

\kbvsechead{(೨)\hspace{0.45em}ಅನುನಾಸಿಕ ಸಂಧಿ :}

\indent ಈ ಉದಾಹರಣೆಗಳನ್ನು ಗಮನಿಸಿ—\par

\indent ವಾಕ್‌ + ಮಯ = ವಾಙ್‌ + ಮಯ = ವಾಙ್ಮಯ (ಕಕಾರಕ್ಕೆ ಮಕಾರ ಪರವಾಗಿದ್ದು ಙ ಕಾರ ಆದೇಶವಾಗಿದೆ)\par
